\documentclass[12pt,a4paper]{article}

% ===== Encodage et langue =====
\usepackage[utf8]{inputenc}
\usepackage[T1]{fontenc}
\usepackage[french]{babel}

% ===== Mise en page =====
\usepackage[margin=2.5cm, headheight=15pt]{geometry}
\usepackage{fancyhdr}
\usepackage{lastpage}

% ===== Couleurs =====
\usepackage{xcolor}
\definecolor{codegreen}{rgb}{0,0.5,0}
\definecolor{codegray}{rgb}{0.5,0.5,0.5}
\definecolor{codepurple}{rgb}{0.58,0,0.82}
\definecolor{backcolour}{rgb}{0.97,0.97,0.97}
\definecolor{ruleblue}{RGB}{0,102,204}
\definecolor{importantred}{RGB}{204,0,0}
\definecolor{tipgreen}{RGB}{0,153,51}
\definecolor{exampleorange}{RGB}{255,140,0}

% ===== Code source =====
\usepackage{listings}
\lstset{
    language=C,
    backgroundcolor=\color{backcolour},
    commentstyle=\color{codegreen},
    keywordstyle=\color{blue},
    numberstyle=\tiny\color{codegray},
    stringstyle=\color{codepurple},
    basicstyle=\ttfamily\small,
    breakatwhitespace=false,
    breaklines=true,
    captionpos=b,
    keepspaces=true,
    numbers=left,
    numbersep=5pt,
    showspaces=false,
    showstringspaces=false,
    showtabs=false,
    tabsize=2,
    frame=single,
    framerule=0.5pt,
    xleftmargin=15pt,
    framexleftmargin=15pt,
    literate={é}{{\'e}}1 {è}{{\`e}}1 {ê}{{\^e}}1 {ë}{{\"e}}1
             {à}{{\`a}}1 {â}{{\^a}}1
             {ù}{{\`u}}1 {û}{{\^u}}1
             {î}{{\^i}}1 {ï}{{\"i}}1
             {ô}{{\^o}}1
             {ç}{{\c{c}}}1
}

% Style C++
\lstdefinestyle{cpp}{
    language=C++,
    backgroundcolor=\color{backcolour},
    commentstyle=\color{codegreen},
    keywordstyle=\color{blue},
    numberstyle=\tiny\color{codegray},
    stringstyle=\color{codepurple},
    basicstyle=\ttfamily\small,
    breakatwhitespace=false,
    breaklines=true,
    captionpos=b,
    keepspaces=true,
    numbers=left,
    numbersep=5pt,
    showspaces=false,
    showstringspaces=false,
    showtabs=false,
    tabsize=2,
    frame=single,
    framerule=0.5pt,
    xleftmargin=15pt,
    framexleftmargin=15pt,
    literate={é}{{\'e}}1 {è}{{\`e}}1 {ê}{{\^e}}1 {ë}{{\"e}}1
             {à}{{\`a}}1 {â}{{\^a}}1
             {ù}{{\`u}}1 {û}{{\^u}}1
             {î}{{\^i}}1 {ï}{{\"i}}1
             {ô}{{\^o}}1
             {ç}{{\c{c}}}1
}

% Style pour bash
\lstdefinestyle{bash}{
    language=bash,
    backgroundcolor=\color{backcolour},
    basicstyle=\ttfamily\small,
    breaklines=true,
    frame=single,
    numbers=left,
    numberstyle=\tiny\color{codegray},
    xleftmargin=15pt,
    framexleftmargin=15pt,
    commentstyle=\color{codegreen},
    keywordstyle=\color{blue},
    morekeywords={gcc,g++,valgrind},
}

% ===== Boites colorées =====
\usepackage[most]{tcolorbox}

\newtcolorbox{regle}{
    colback=ruleblue!5,
    colframe=ruleblue,
    fonttitle=\bfseries,
    title=R\`egle,
    sharp corners,
    boxrule=1pt,
    left=5pt, right=5pt, top=5pt, bottom=5pt,
}

\newtcolorbox{important}{
    colback=importantred!5,
    colframe=importantred,
    fonttitle=\bfseries,
    title=Important,
    sharp corners,
    boxrule=1pt,
    left=5pt, right=5pt, top=5pt, bottom=5pt,
}

\newtcolorbox{astuce}{
    colback=tipgreen!5,
    colframe=tipgreen,
    fonttitle=\bfseries,
    title=Astuce,
    sharp corners,
    boxrule=1pt,
    left=5pt, right=5pt, top=5pt, bottom=5pt,
}

\newtcolorbox{exemple}{
    colback=exampleorange!5,
    colframe=exampleorange,
    fonttitle=\bfseries,
    title=Exemple,
    sharp corners,
    boxrule=1pt,
    left=5pt, right=5pt, top=5pt, bottom=5pt,
}

% ===== Mathématiques =====
\usepackage{amsmath}

% ===== Tableaux =====
\usepackage{array}
\usepackage{booktabs}

% ===== Liens =====
\usepackage{hyperref}
\hypersetup{
    colorlinks=true,
    linkcolor=ruleblue,
    urlcolor=ruleblue,
    citecolor=ruleblue,
    bookmarks=true,
    bookmarksopen=true,
    bookmarksnumbered=true,
    linktoc=all,
    pdftitle={Cours complet C et C++},
    pdfauthor={Arthur},
    pdfsubject={L1 Informatique -- UAPV},
}

% ===== En-tête et pied de page =====
\pagestyle{fancy}
\fancyhf{}
\fancyhead[L]{Arthur -- L1 Informatique -- Cours C / C++}
\fancyfoot[C]{\thepage}
\renewcommand{\headrulewidth}{0.4pt}
\renewcommand{\footrulewidth}{0pt}

% ===== Titre =====
\title{\Huge\textbf{Cours complet C \& C++}\\[0.5cm]
\Large De l'algorithmique aux objets\\[0.3cm]
\large S1--S2 -- Licence Informatique -- UAPV}
\author{Arthur}
\date{2025 / 2026}

% ============================================================
\begin{document}

\maketitle
\thispagestyle{empty}

\newpage
\tableofcontents
\newpage

% %%%%%%%%%%%%%%%%%%%%%%%%%%%%%%%%%%%%%%%%%%%%%%%%%%%%%%%%%%%%
%                        PARTIE I : C
% %%%%%%%%%%%%%%%%%%%%%%%%%%%%%%%%%%%%%%%%%%%%%%%%%%%%%%%%%%%%

\part{Bases de la programmation en C}

% ============================================================
\section{Algorithmes, variables, conditions et boucles}

\subsection{C'est quoi un algorithme ?}

Un algorithme est une \textbf{suite finie d'instructions}, exécutées dans l'ordre, qui prend des entrées et produit un résultat.
Un bon algorithme est \textbf{peu complexe}, \textbf{rapide} et \textbf{économe en ressources}.

\begin{exemple}
Deux façons de représenter un algorithme :
\begin{itemize}
    \item \textbf{Pseudo-code} : description textuelle proche du langage humain
    \item \textbf{Algorigramme} : représentation graphique avec des boîtes et des flèches
\end{itemize}
\end{exemple}

\subsection{C vs Python — bas niveau vs haut niveau}

\begin{table}[h!]
\centering
\caption{Comparaison C / Python}
\begin{tabular}{lll}
\toprule
 & \textbf{C / C++} & \textbf{Python} \\
\midrule
Niveau & Bas (proche machine) & Haut (proche humain) \\
Gestion mémoire & Manuelle & Automatique \\
Compilation & Obligatoire (gcc / g++) & Interprété \\
Performance & Très rapide & Plus lent \\
Portabilité & Bonne & Excellente \\
\bottomrule
\end{tabular}
\end{table}

Le compilateur (\texttt{gcc} pour le C, \texttt{g++} pour le C++) transforme le code source en langage machine, effectue des analyses lexicale, syntaxique et sémantique, signale les erreurs et émet des avertissements.

\begin{lstlisting}[style=bash, caption={Compilation et exécution d'un programme C}]
# Compilation
gcc programme.c -o programme

# Exécution
./programme
\end{lstlisting}

\subsection{Structure minimale d'un programme C}

\begin{lstlisting}[caption={Hello World en C}]
#include <stdio.h>

int main() {
    printf("Hello World\n");
    return 0;
}
\end{lstlisting}

\begin{itemize}
    \item \texttt{\#include <stdio.h>} : inclut la bibliothèque d'entrées/sorties standard
    \item \texttt{int main()} : point d'entrée, retourne 0 en cas de succès
    \item \texttt{printf} : affiche du texte dans le terminal
    \item Chaque instruction se termine par \texttt{;}
\end{itemize}

\subsection{Les variables}

Une variable possède : un \textbf{nom}, un \textbf{type}, une \textbf{valeur} et une \textbf{adresse} en mémoire.

\subsubsection{Déclaration et initialisation}

\begin{lstlisting}[caption={Déclaration et initialisation de variables}]
// Déclaration
int age;
double taille;
bool vivant;

// Initialisation (affectation)
age = 20;
taille = 1.78;
vivant = true; // nécessite #include <stdbool.h> en C

// Déclaration + initialisation en une ligne
int score = 100;
double prix_ht = 9.99;
\end{lstlisting}

\subsubsection{Types principaux}

\begin{table}[h!]
\centering
\caption{Types de base en C}
\begin{tabular}{llll}
\toprule
\textbf{Type} & \textbf{Taille} & \textbf{Description} & \textbf{Exemple} \\
\midrule
\texttt{int}    & 32 bits & Entier           & \texttt{int x = 42;} \\
\texttt{double} & 64 bits & Réel double prec. & \texttt{double pi = 3.14;} \\
\texttt{float}  & 32 bits & Réel simple prec. & \texttt{float f = 2.5f;} \\
\texttt{char}   & 8 bits  & Caractère        & \texttt{char c = 'A';} \\
\texttt{bool}   & 8 bits  & Booléen (C99)    & \texttt{bool ok = true;} \\
\bottomrule
\end{tabular}
\end{table}

\subsubsection{Convention de nommage (snake\_case)}

\begin{itemize}
    \item En anglais, en minuscules
    \item Commence par une lettre de l'alphabet
    \item Séparateur : \texttt{\_} (ex : \texttt{prix\_final}, \texttt{length\_array})
    \item Pas de mots-clés réservés (\texttt{int}, \texttt{return}, etc.)
\end{itemize}

\subsubsection{Entrée / sortie : printf et scanf}

\begin{lstlisting}[caption={Entrées et sorties en C}]
#include <stdio.h>

int main() {
    int age;
    double taille;

    // Affichage
    printf("Bonjour !\n");
    printf("Score : %d\n", 42);         // %d = int
    printf("Pi = %f\n", 3.14159);       // %f = double
    printf("Lettre : %c\n", 'A');       // %c = char

    // Saisie utilisateur
    printf("Entrez votre age : ");
    scanf("%d", &age);                  // & obligatoire !

    printf("Entrez votre taille : ");
    scanf("%lf", &taille);              // %lf pour double

    printf("Vous avez %d ans et mesurez %f m\n", age, taille);
    return 0;
}
\end{lstlisting}

\begin{regle}
Correspondance format / type pour \texttt{printf} et \texttt{scanf} :
\begin{itemize}
    \item \texttt{\%d} -- \texttt{int}
    \item \texttt{\%f} -- \texttt{double} (printf) / \texttt{\%lf} -- \texttt{double} (scanf)
    \item \texttt{\%c} -- \texttt{char}
    \item \texttt{\%p} -- adresse mémoire (pointeur)
    \item Toujours mettre \texttt{\&} devant la variable dans \texttt{scanf} !
\end{itemize}
\end{regle}

\subsubsection{Portée des variables}

Une variable n'existe que dans le \textbf{bloc} (\texttt{\{\}}) où elle est déclarée, et dans les blocs imbriqués. Dès que l'on sort du bloc, la variable est détruite.

\begin{lstlisting}[caption={Portée locale}]
int main() {
    int a = 3;

    if (a < 5) {
        int c = 5;   // c existe seulement ici
    }

    // printf("%d\n", c); // ERREUR : c n'existe plus
    return 0;
}
\end{lstlisting}

\subsection{Les conditions}

\begin{lstlisting}[caption={Conditions if / else if / else}]
int doigt = 2;

if (doigt == 1) {
    printf("THUMB UP !\n");
} else if (doigt == 2) {
    printf("ON MONTRE PAS DU DOIGT !\n");
} else if (doigt == 3) {
    printf("SOIS POLI !\n");
} else if (doigt == 4 || doigt == 5) {
    printf("ILS SERVENT A RIEN\n");
} else {
    printf("C'est quoi cette tentacule ?!\n");
}
\end{lstlisting}

\subsubsection{Opérateurs relationnels}

\begin{table}[h!]
\centering
\begin{tabular}{ll}
\toprule
\textbf{Opérateur} & \textbf{Signification} \\
\midrule
\texttt{==} & Égal \\
\texttt{!=} & Différent \\
\texttt{<}  & Inférieur \\
\texttt{<=} & Inférieur ou égal \\
\texttt{>}  & Supérieur \\
\texttt{>=} & Supérieur ou égal \\
\bottomrule
\end{tabular}
\end{table}

\subsubsection{Opérateurs logiques}

\begin{table}[h!]
\centering
\begin{tabular}{ll}
\toprule
\textbf{Opérateur} & \textbf{Signification} \\
\midrule
\texttt{\&\&} & ET logique \\
\texttt{||} & OU logique \\
\texttt{!} & NON (inverse la valeur booléenne) \\
\bottomrule
\end{tabular}
\end{table}

\begin{lstlisting}[caption={Exemples d'opérateurs logiques}]
if (a > 0 && b > 0)              // a ET b positifs
if (a < 0 || b > 0)              // a négatif OU b positif
if (a > 0 && (b > 0 || c > 0))  // a positif ET (b OU c positifs)
if (a > 0 && !(b > 0 || c > 0)) // a positif ET NI b NI c positifs
\end{lstlisting}

\subsection{Les boucles}

\subsubsection{Boucle for — nombre d'itérations connu}

\begin{lstlisting}[caption={Boucle for}]
// for (initialisation ; condition de continuation ; pas)
for (int i = 0; i < 10; i = i + 1) {
    printf("%d\n", i);
}
\end{lstlisting}

\subsubsection{Boucle while — nombre d'itérations inconnu}

\begin{lstlisting}[caption={Boucle while}]
int i = 0;
while (i < 10) {
    i = chiffreAuHasard();
}
\end{lstlisting}

\begin{regle}
\begin{itemize}
    \item \textbf{for} : quand on connaît à l'avance le nombre d'itérations.
    \item \textbf{while} : quand on répète jusqu'à ce qu'une condition soit satisfaite.
\end{itemize}
\end{regle}

% ============================================================
\section{Les fonctions}

\subsection{Qu'est-ce qu'une fonction ?}

Une fonction est une \textbf{boîte qui encapsule un algorithme} : elle prend des entrées (paramètres), effectue un traitement et retourne un résultat. Elle permet de ne pas répéter du code, de le rendre lisible et de le réutiliser.

\subsection{Déclaration et définition}

\begin{lstlisting}[caption={Anatomie d'une fonction}]
// prototype : type_retour nom(type param1, type param2)
int additionner(int nombre1, int nombre2)
{
    int result = nombre1 + nombre2;
    return result;
}

int main() {
    int a = 12, b = 18;
    int somme = additionner(a, b); // appel
    printf("Somme : %d\n", somme);
    return 0;
}
\end{lstlisting}

\begin{regle}
Convention de nommage des fonctions (camelCase) :
\begin{itemize}
    \item Commence par un verbe, en minuscule : \texttt{calculerMoyenne}, \texttt{afficherScore}
    \item Chaque nouveau mot commence par une majuscule
    \item Pas de symboles spéciaux
\end{itemize}
\end{regle}

\subsection{Fonctions sans retour : void}

\begin{lstlisting}[caption={Fonction void}]
void afficherSalutation(int age) {
    if (age >= 18) {
        printf("Bienvenue !\n");
    } else {
        printf("Interdit aux mineurs.\n");
    }
    // pas de return nécessaire
}
\end{lstlisting}

\subsection{Passage par valeur (copie)}

En C, les paramètres sont \textbf{toujours passés par valeur} : la fonction reçoit une \textbf{copie} de la variable. Modifier le paramètre dans la fonction ne modifie pas la variable originale.

\begin{lstlisting}[caption={Passage par valeur — la variable originale n'est pas modifiée}]
void addOne(int a) {
    a = a + 1; // modifie la COPIE
}

int main() {
    int a = 15;
    addOne(a);
    printf("%d\n", a); // affiche 15, pas 16 !
    return 0;
}
\end{lstlisting}

\begin{important}
En Python, le passage par défaut est le \textbf{passage par référence} pour les objets mutables. C'est l'inverse du C !
\end{important}

\subsection{Retourner un résultat}

\texttt{return} renvoie une valeur \textbf{et quitte immédiatement} la fonction.

\begin{lstlisting}[caption={Retour de valeur et return early}]
int addOne(int number) {
    number = number + 1;
    return number; // retourne la valeur modifiée
}

// return quitte la fonction immédiatement :
bool isAboveTen(int number) {
    if (number > 10) {
        return true;  // on sort ici si vrai
    }
    return false;     // équivalent à l'else
}
\end{lstlisting}

\subsection{La pile (stack) et les fonctions}

À chaque appel de fonction, une nouvelle \textbf{frame} est empilée sur la pile. Les variables locales de cette fonction y vivent. À la sortie de la fonction, la frame est dépilée automatiquement et les variables sont détruites.

\begin{lstlisting}[caption={Pile d'exécution}]
void c() { int z = 3; }
void b() { int y = 2; }
void a() { int x = 1; b(); }

int main() {
    int w = 3;
    if (w == 3) {
        int v = 5; // portée locale au if
        a();
    }
    c();
    return 0;
}
// Pile au moment de b() :
//   b()  : y
//   a()  : x
//   if   : v
//   main : w
\end{lstlisting}

\subsection{Fichiers .h et .c}

Pour organiser un projet, on sépare les \textbf{déclarations} (fichier \texttt{.h}) des \textbf{définitions} (fichier \texttt{.c}).

\begin{lstlisting}[caption={mathematicOperations.h}]
#ifndef MATHOPERATION_H
#define MATHOPERATION_H

int addOne(int a);
double sum(double a, double b);
int power(int a, int pow);

#endif
\end{lstlisting}

\begin{lstlisting}[caption={mathematicOperations.c}]
#include "mathematicOperations.h"

int addOne(int a) {
    a = a + 1;
    return a;
}

double sum(double a, double b) {
    double result = a + b;
    return result;
}
\end{lstlisting}

\begin{lstlisting}[caption={main.c}]
#include <stdio.h>
#include "mathematicOperations.h"

int main() {
    int result = addOne(12);
    printf("Résultat : %d\n", result);
    return 0;
}
\end{lstlisting}

% ============================================================
\section{Les pointeurs}

\subsection{Rappel : la mémoire}

Toute variable possède une \textbf{adresse} (numéro de case mémoire) et un \textbf{type} (qui dit combien de bits lire). Un pointeur stocke justement cette adresse.

\begin{table}[h!]
\centering
\caption{Exemple de représentation mémoire}
\begin{tabular}{lll}
\toprule
\textbf{Variable} & \textbf{Adresse} & \textbf{Valeur} \\
\midrule
\texttt{a}    & \texttt{0x7ffcc0bee04c} & \texttt{3} \\
\texttt{pAge} & \texttt{0x9ffdce1f45ee} & \texttt{0x7ffcc0bee04c} \\
\bottomrule
\end{tabular}
\end{table}

\subsection{Déclarer et utiliser un pointeur}

\begin{lstlisting}[caption={Déclaration, & et *}]
int a = 3;

// Déclaration d'un pointeur (convention : nom commence par p)
int *pAge;

// Récupérer l'adresse de a et la stocker dans pAge
pAge = &a;

// Afficher l'adresse et la valeur pointée
printf("%d\n", a);      // 3
printf("%p\n", pAge);   // 0x7ffcc...
printf("%d\n", *pAge);  // 3  (déréférencement)

// Modifier la valeur via le pointeur
*pAge = 6;
printf("%d\n", a);      // 6  (a est modifié !)
printf("%d\n", *pAge);  // 6
\end{lstlisting}

\begin{regle}
\begin{itemize}
    \item \texttt{\&} devant une variable : donne son \textbf{adresse}
    \item \texttt{*} à la \textbf{déclaration} : indique que la variable est un pointeur
    \item \texttt{*} à \textbf{l'utilisation} : accède à la valeur pointée (déréférencement)
\end{itemize}
\end{regle}

\subsection{Erreurs courantes}

\begin{table}[h!]
\centering
\caption{Erreurs fréquentes avec les pointeurs}
\begin{tabular}{lp{8cm}}
\toprule
\textbf{Erreur} & \textbf{Cause} \\
\midrule
Null Pointer Exception & Pointeur non initialisé, mis à \texttt{NULL} ou désalloué \\
Invalid Pointer Error  & Adresse modifiée par erreur, ou espace mémoire désalloué \\
\bottomrule
\end{tabular}
\end{table}

\begin{important}
Toujours initialiser un pointeur ! Si vous ne pointez sur rien encore, mettez-le à \texttt{NULL} :
\begin{lstlisting}
int *pAge = NULL;
\end{lstlisting}
\end{important}

\subsection{Pointeurs comme paramètres (in-out)}

Le passage par pointeur permet de \textbf{modifier la variable originale} depuis une fonction — ce que le passage par valeur ne permet pas.

\begin{lstlisting}[caption={Passage par pointeur — modifie l'original}]
void addOne(int *pBase) {
    *pBase = *pBase + 1;
}

int main() {
    int a = 3;
    addOne(&a);         // on passe l'ADRESSE de a
    printf("%d\n", a);  // affiche 4 !
    return 0;
}
\end{lstlisting}

\begin{astuce}
En C, si une fonction doit retourner \textbf{plusieurs résultats}, on passe des pointeurs en paramètres (paramètres in-out). En C++, on utilisera les références pour simplifier cette syntaxe.
\end{astuce}

\subsection{Pointeur de pointeur}

\begin{lstlisting}[caption={Double pointeur}]
int a = 42;
int *pB = &a;
int **pPC = &pB;  // pointeur vers un pointeur

printf("%p\n",  pPC);  // adresse de pB
printf("%p\n", *pPC);  // adresse de a
printf("%d\n", **pPC); // valeur de a : 42
\end{lstlisting}

% ============================================================
\section{Les tableaux}

\subsection{Motivation}

Plutôt que 300 variables séparées pour stocker les notes de 300 étudiants, un tableau permet de tout regrouper en une seule variable et de tout traiter avec une boucle.

\subsection{Déclaration et initialisation}

\begin{lstlisting}[caption={Déclaration et initialisation de tableaux}]
// Déclaration (taille fixe à la compilation)
int tab[300];

// Initialisation totale
int tab5[5] = {3, 5, 8, 12, 8};
double tabD[3] = {3.5, 8.2, 7.3};
bool tabB[4] = {0, 1, 1, 1};

// Initialisation partielle (le reste vaut 0)
int tab2[5] = {3, 5, 8};  // -> {3, 5, 8, 0, 0}

// Initialisation après déclaration (case par case)
int tab3[5];
tab3[0] = 0;
tab3[1] = 2;
// ou avec une boucle :
for (int i = 0; i < 5; i++) {
    tab3[i] = i * 2;
}
\end{lstlisting}

\subsection{Accès aux données}

La numérotation commence à \textbf{0}. \texttt{tab[i]} accède à la case d'indice \texttt{i} (décalage de \texttt{i} depuis le début du tableau).

\begin{lstlisting}[caption={Accès et modification}]
int tab[4] = {3, 8, 56, 158};

printf("%d\n", tab[0]);  // 3   (1ère case)
printf("%d\n", tab[1]);  // 8   (2ème case)
printf("%d\n", tab[3]);  // 158 (4ème case)

tab[3] = 14;             // modifie la 4ème case
\end{lstlisting}

\subsection{Tableau et pointeur}

Un tableau se comporte comme un \textbf{pointeur constant} vers sa première case. La variable \texttt{tab} contient l'adresse du premier élément.

\begin{lstlisting}[caption={Équivalence tableau / pointeur}]
int tab[4] = {3, 8, 56, 158};

printf("%p\n",  tab);          // adresse de tab[0]
printf("%d\n", *tab);          // 3  (équivalent à tab[0])
printf("%d\n", *(tab + 1));    // 8  (équivalent à tab[1])
\end{lstlisting}

\subsection{Taille statique vs taille réelle}

\begin{lstlisting}[caption={Gérer taille max et taille réelle}]
#define MAX_ETUDIANTS 2000

int etudiants[MAX_ETUDIANTS];
int tailleMax = MAX_ETUDIANTS;
int tailleReelle;

printf("Nombre d'étudiants cette année : ");
scanf("%d", &tailleReelle);

for (int i = 0; i < tailleReelle; i++) {
    scanf("%d", &etudiants[i]);
}
\end{lstlisting}

\begin{important}
La taille d'un tableau statique est fixée à la compilation. On ne peut pas écrire \texttt{int tab[n];} si \texttt{n} est saisi par l'utilisateur (avant C99). La solution propre : réserver un max et stocker la taille réelle séparément.
\end{important}

\subsection{Tableaux et fonctions}

Un tableau passé à une fonction l'est comme un pointeur : les modifications dans la fonction impactent le tableau original.

\begin{lstlisting}[caption={Passage de tableau en paramètre}]
// Toujours passer la taille avec le tableau !
void affiche(int tab[], int taille) {
    for (int i = 0; i < taille; i++) {
        printf("%d\n", tab[i]);
    }
}

void saisieTableau(int tableau[], int taille) {
    for (int i = 0; i < taille; i++) {
        scanf("%d", &tableau[i]);
    }
}

int main() {
    int tab[4] = {2, 5, 8, 9};
    affiche(tab, 4);
    return 0;
}
\end{lstlisting}

\subsection{Tableaux multidimensionnels}

\begin{lstlisting}[caption={Tableau 2D}]
// Déclaration : int tab[lignes][colonnes]
int tab[3][4];

// Initialisation
int mat[3][4] = { {0, 1, 2, 3},
                  {4, 5, 6, 7},
                  {8, 9, 10, 11} };

// Parcours avec boucles imbriquées
for (int i = 0; i < 3; i++) {
    for (int j = 0; j < 4; j++) {
        printf("%d ", mat[i][j]);
    }
    printf("\n");
}

// Passage en fonction
void saisieTab2D(int tableau[][4], int tailleX, int tailleY);
\end{lstlisting}

% ============================================================
\section{Les structures}

\subsection{Motivation}

Quand plusieurs variables décrivent la même entité (nom, date de naissance, numéro de sécu d'une personne), il devient difficile de les gérer séparément. Les structures permettent de les \textbf{regrouper en un seul type}.

\subsection{Définir une structure}

Les structures se définissent dans un fichier \texttt{.h} portant leur nom.

\begin{lstlisting}[caption={Tamagotchi.h}]
#ifndef TAMAGOTCHI_H
#define TAMAGOTCHI_H

typedef struct {
    int   sante;
    int   energie;
    double faim;
    double moral;
} Tamagotchi;

#endif
\end{lstlisting}

\begin{regle}
Règles de nommage d'une structure :
\begin{itemize}
    \item Commence par une majuscule, c'est un \textbf{nom} (pas un verbe)
    \item Pas de \texttt{\_} ni de caractères spéciaux
    \item N'oubliez pas le \texttt{;} après la fermeture \texttt{\}}
\end{itemize}
\end{regle}

\subsection{Utiliser une structure}

\begin{lstlisting}[caption={Déclaration, initialisation et accès}]
#include "Tamagotchi.h"
#include <stdio.h>

int main() {
    // Déclaration
    Tamagotchi poupi;

    // Accès et modification : NomInstance.NomChamp
    poupi.sante  = 12;
    poupi.energie = 80;

    printf("Santé : %d\n", poupi.sante);

    // Initialisation séquentielle à la déclaration
    Tamagotchi riri = {100, 100, 0.0, 100.0};

    // Initialisation sélective (le reste vaut 0)
    Tamagotchi loulou = {.energie = 100, .moral = 100.0};

    // Tableau de structures
    Tamagotchi tamadex[151];
    tamadex[0].sante = 24;

    return 0;
}
\end{lstlisting}

\subsection{Fonctions et structures}

\begin{lstlisting}[caption={Tamagotchi.c — associer des fonctions à une structure}]
#include "Tamagotchi.h"
#include <stdio.h>

// Passage par copie : les modifications n'impactent pas l'original
void afficherSante(Tamagotchi tama) {
    printf("Santé : %d\n", tama.sante);
}

// Retourner une structure modifiée
Tamagotchi setSante(Tamagotchi tama, int sante) {
    if      (sante < 0)   tama.sante = 0;
    else if (sante > 100) tama.sante = 100;
    else                  tama.sante = sante;
    return tama;
}

// Passage par pointeur : modifie l'original directement
void setSantePtr(Tamagotchi *pTama, int sante) {
    if      (sante < 0)   (*pTama).sante = 0;
    else if (sante > 100) (*pTama).sante = 100;
    else                  (*pTama).sante = sante;
}
\end{lstlisting}

\begin{astuce}
\texttt{(*pTama).sante} et \texttt{pTama->sante} sont équivalents. L'opérateur \texttt{->} est une notation raccourcie très utilisée en C pour accéder à un champ via un pointeur.
\end{astuce}

\begin{lstlisting}[caption={Utilisation dans main.c}]
int main() {
    Tamagotchi poupi = {100, 100, 0.0, 100.0};

    // Passage par copie + récupération du résultat
    poupi = setSante(poupi, 150);   // sera clampé à 100
    afficherSante(poupi);

    // Passage par pointeur
    setSantePtr(&poupi, -5);        // sera clampé à 0
    printf("Santé : %d\n", poupi.sante);

    return 0;
}
\end{lstlisting}

% ============================================================
\section{Allocation dynamique}

\subsection{Pourquoi allouer dynamiquement ?}

Les variables locales vivent sur la \textbf{pile} et sont détruites en sortant de leur bloc. L'\textbf{allocation dynamique} permet de créer de la mémoire sur le \textbf{tas} (heap), dont la durée de vie n'est pas liée à la portée des variables.

\begin{table}[h!]
\centering
\caption{Pile (Stack) vs Tas (Heap)}
\begin{tabular}{lll}
\toprule
 & \textbf{Pile (Stack)} & \textbf{Tas (Heap)} \\
\midrule
Gestion & Automatique & Manuelle \\
Durée de vie & Liée au bloc & Jusqu'à \texttt{free()} \\
Taille & Limitée ($\sim$8 Mo) & Quasi-illimitée \\
Vitesse & Rapide & Plus lent \\
\bottomrule
\end{tabular}
\end{table}

\subsection{malloc — allouer de la mémoire}

\begin{lstlisting}[caption={Allocation avec malloc}]
#include <stdlib.h>

// Allouer un entier sur le tas
int *pVar = malloc(sizeof(int));
*pVar = 3;
printf("%d\n", *pVar);  // 3

// sizeof retourne la taille en octets d'un type
// sizeof(int) -> 4, sizeof(double) -> 8, etc.
\end{lstlisting}

\begin{important}
Toujours stocker l'adresse renvoyée par \texttt{malloc} dans un pointeur ! Sinon l'espace est alloué mais introuvable.
\end{important}

\subsection{free — libérer la mémoire}

\begin{lstlisting}[caption={Libération avec free}]
int *pAge = malloc(sizeof(int));
*pAge = 3;

// Libération
free(pAge);
pAge = NULL;  // Bonne pratique : mettre à NULL après free
\end{lstlisting}

\begin{important}
\textbf{Fuite mémoire} : si on perd l'adresse avant d'avoir fait \texttt{free}, la mémoire reste allouée jusqu'à la fin du programme. Chaque \texttt{malloc} doit avoir son \texttt{free}.
\end{important}

\subsection{Tableau alloué dynamiquement}

\begin{lstlisting}[caption={Tableau de taille choisie à l'exécution}]
int n;
printf("Taille du tableau : ");
scanf("%d", &n);

// Allouer n entiers contigus
int *arr = malloc(n * sizeof(int));

// Utilisation identique à un tableau classique
arr[5] = 3;
arr[0] = 10;

// Libération
free(arr);
arr = NULL;
\end{lstlisting}

\subsection{malloc et structures}

\begin{lstlisting}[caption={Allocation dynamique d'une structure}]
// Structure avec un tableau interne
typedef struct {
    int    sante;
    double *stats;
} Tamagotchi;

int main() {
    // Allocation de la structure
    Tamagotchi *t = malloc(sizeof(Tamagotchi));
    (*t).sante = 100;

    // Allocation du tableau interne
    int n = 5;
    (*t).stats = malloc(n * sizeof(double));
    (*t).stats[0] = 9.5;

    // Libération (ordre inverse de l'allocation)
    free((*t).stats);
    free(t);
    t = NULL;

    return 0;
}
\end{lstlisting}

\begin{regle}
Ordre de libération : toujours libérer les allocations \textbf{imbriquées en premier}, puis la structure principale. L'ordre inverse de l'allocation.
\end{regle}

% %%%%%%%%%%%%%%%%%%%%%%%%%%%%%%%%%%%%%%%%%%%%%%%%%%%%%%%%%%%%
%                       PARTIE II : C++
% %%%%%%%%%%%%%%%%%%%%%%%%%%%%%%%%%%%%%%%%%%%%%%%%%%%%%%%%%%%%

\part{Programmation en C++}

% ============================================================
\section{Introduction au C++}

\subsection{C++ par rapport au C}

Le C++ est une \textbf{extension du C} qui ajoute la programmation orientée objet, les références, la surcharge de fonctions, les templates, et une bibliothèque standard riche. Tout code C valide est (presque) valide en C++.

\begin{table}[h!]
\centering
\caption{Principales différences C / C++}
\begin{tabular}{lll}
\toprule
\textbf{Fonctionnalité} & \textbf{C} & \textbf{C++} \\
\midrule
Entrées/sorties & \texttt{printf / scanf} & \texttt{cout / cin} \\
Allocation dynamique & \texttt{malloc / free} & \texttt{new / delete} \\
Références & Non & Oui (\texttt{\&}) \\
Classes et objets & Non (struct seulement) & Oui \\
Surcharge de fonctions & Non & Oui \\
Passage par référence & Par pointeur & Par référence (\texttt{\&}) \\
\bottomrule
\end{tabular}
\end{table}

\subsection{Caractéristiques du langage}

Le C++ est un langage \textbf{compilé}, c'est-à-dire qu'avant d'être exécuté, le code source doit être traduit en langage machine par un compilateur.

\begin{exemple}
Processus de compilation et exécution :
\begin{enumerate}
    \item Écriture du code dans \texttt{programme.cpp}
    \item Compilation : \texttt{g++ programme.cpp -o programme}
    \item Exécution : \texttt{./programme}
\end{enumerate}
\end{exemple}

\subsection{Structure minimale d'un programme}

\begin{lstlisting}[style=cpp, caption={Programme minimal en C++}]
#include <iostream>
using namespace std;

int main() {
    cout << "Hello World!" << endl;
    return 0;
}
\end{lstlisting}

\begin{itemize}
    \item \texttt{\#include <iostream>} : inclut la bibliothèque d'entrées/sorties
    \item \texttt{using namespace std;} : permet d'utiliser \texttt{cout} et \texttt{cin} sans préfixe
    \item \texttt{int main()} : fonction principale, point d'entrée du programme
\end{itemize}

\subsection{Entrées et sorties}

En C++, on utilise \texttt{cout} et \texttt{cin} à la place de \texttt{printf} et \texttt{scanf}.

\begin{lstlisting}[style=cpp, caption={Utilisation de cout et cin}]
int age;
string nom;

// Affichage
cout << "Entrez votre nom : ";
cin >> nom;

cout << "Entrez votre age : ";
cin >> age;

cout << "Bonjour " << nom << ", vous avez " << age << " ans." << endl;
\end{lstlisting}

\begin{regle}
Opérateurs de flux :
\begin{itemize}
    \item \texttt{<<} : pour envoyer vers la sortie (\texttt{cout})
    \item \texttt{>>} : pour récupérer depuis l'entrée (\texttt{cin})
    \item \texttt{endl} : retour à la ligne
\end{itemize}
\end{regle}

% ============================================================
\section{Les types de base et les tableaux}

\subsection{Types de données}

\begin{table}[h!]
\centering
\caption{Types de données en C++}
\begin{tabular}{lll}
\toprule
\textbf{Type} & \textbf{Description} & \textbf{Exemple} \\
\midrule
\texttt{int} & Entier & \texttt{int x = 42;} \\
\texttt{double} & Réel double précision & \texttt{double pi = 3.14;} \\
\texttt{float} & Réel simple précision & \texttt{float f = 2.5f;} \\
\texttt{char} & Caractère & \texttt{char c = 'A';} \\
\texttt{bool} & Booléen & \texttt{bool ok = true;} \\
\texttt{string} & Chaîne de caractères & \texttt{string s = "hello";} \\
\bottomrule
\end{tabular}
\end{table}

\subsection{Déclaration de tableaux}

\begin{lstlisting}[style=cpp, caption={Tableaux statiques}]
// Tableau a une dimension
int tab[10];       // Tableau de 10 entiers
tab[0] = 5;        // Premier element (indice 0)
tab[9] = 12;       // Dernier element (indice 9)

// Tableau a deux dimensions
double matrice[3][4]; // 3 lignes, 4 colonnes
matrice[0][0] = 1.5;

// Initialisation lors de la declaration
int nombres[] = {10, 20, 30, 40, 50};
\end{lstlisting}

\begin{important}
Les tableaux statiques ont une taille fixe définie à la compilation. Cette taille ne peut pas être modifiée pendant l'exécution du programme.
\end{important}

% ============================================================
\section{Pointeurs et références}

\subsection{Les pointeurs}

Un pointeur est une variable qui contient l'\textbf{adresse mémoire} d'une autre variable.

\begin{lstlisting}[style=cpp, caption={Déclaration et utilisation de pointeurs}]
int x = 42;
int* p;          // p est un pointeur sur int

p = &x;          // p contient l'adresse de x

cout << x;       // Affiche : 42
cout << &x;      // Affiche : 0x7fff5a3b524 (adresse de x)
cout << p;       // Affiche : 0x7fff5a3b524 (meme adresse)
cout << *p;      // Affiche : 42 (valeur pointee)
\end{lstlisting}

\begin{regle}
Les deux opérateurs clés :
\begin{itemize}
    \item \texttt{\&} : opérateur d'adresse, donne l'adresse d'une variable
    \item \texttt{*} : opérateur de déréférencement, accède à la valeur pointée
\end{itemize}
\end{regle}

\subsection{Schéma mémoire}

\begin{table}[h!]
\centering
\begin{tabular}{lll}
\toprule
\textbf{Variable} & \textbf{Adresse} & \textbf{Valeur} \\
\midrule
\texttt{x} & \texttt{0x1000} & \texttt{42} \\
\texttt{p} & \texttt{0x2000} & \texttt{0x1000} \\
\bottomrule
\end{tabular}
\end{table}

Le pointeur \texttt{p} (stocké à l'adresse \texttt{0x2000}) contient l'adresse de \texttt{x} (\texttt{0x1000}).

\subsection{Les références}

Une référence est un \textbf{alias} pour une variable existante. C'est la grande nouveauté du C++ par rapport au C.

\begin{lstlisting}[style=cpp, caption={Utilisation des références}]
int x = 10;
int& ref = x;    // ref est une reference sur x

ref = 20;        // Modifie x
cout << x;       // Affiche : 20
\end{lstlisting}

\begin{astuce}
Une référence doit être initialisée lors de sa déclaration et ne peut plus changer de cible par la suite.
\end{astuce}

% ============================================================
\section{Allocation dynamique}

\subsection{La pile (Stack) et le tas (Heap)}

\begin{table}[h!]
\centering
\caption{Comparaison Pile vs Tas}
\begin{tabular}{lll}
\toprule
 & \textbf{Pile (Stack)} & \textbf{Tas (Heap)} \\
\midrule
\textbf{Taille} & Limitée ($\sim$8 Mo) & Quasi-illimitée \\
\textbf{Libération} & Automatique & Manuelle (\texttt{delete}) \\
\textbf{Utilisation} & Variables locales & Allocation dynamique \\
\textbf{Rapidité} & Rapide & Plus lent \\
\bottomrule
\end{tabular}
\end{table}

\subsection{Allocation dans la pile}

\begin{lstlisting}[style=cpp, caption={Variables locales dans la pile}]
void fonction() {
    int x = 5;          // Cree dans la pile
    double tab[10];      // Cree dans la pile
    // ...
} // x et tab sont automatiquement detruits ici
\end{lstlisting}

\subsection{Allocation dans le tas avec new/delete}

En C++, on remplace \texttt{malloc/free} par \texttt{new/delete} qui sont plus sûrs et gèrent les constructeurs.

\begin{lstlisting}[style=cpp, caption={Allocation dynamique avec new}]
// Allocation d'un seul element
int* p = new int;
*p = 42;
delete p;            // Libere la memoire

// Allocation d'un tableau
int n = 100;
int* tableau = new int[n];
tableau[0] = 10;
tableau[99] = 20;
delete[] tableau;    // ATTENTION aux crochets !
\end{lstlisting}

\begin{important}
Règle d'or :
\begin{itemize}
    \item Chaque \texttt{new} doit avoir son \texttt{delete}
    \item Chaque \texttt{new[]} doit avoir son \texttt{delete[]}
    \item Sinon : \textbf{fuite mémoire !}
\end{itemize}
\end{important}

\subsection{Exemple complet}

\begin{lstlisting}[style=cpp, caption={Gestion dynamique d'un tableau}]
#include <iostream>
using namespace std;

int main() {
    int taille;
    cout << "Entrez la taille du tableau : ";
    cin >> taille;

    // Allocation dynamique
    int* tab = new int[taille];

    // Remplissage
    for (int i = 0; i < taille; i++) {
        tab[i] = i * 10;
    }

    // Affichage
    for (int i = 0; i < taille; i++) {
        cout << tab[i] << " ";
    }
    cout << endl;

    // Liberation
    delete[] tab;

    return 0;
}
\end{lstlisting}

% ============================================================
\section{Passage de paramètres}

\subsection{Passage par valeur}

Le paramètre reçoit une \textbf{copie} de la variable.

\begin{lstlisting}[style=cpp, caption={Passage par valeur}]
void modifier(int x) {
    x = 10;          // Modifie la COPIE
}

int main() {
    int a = 5;
    modifier(a);
    cout << a;       // Affiche : 5 (non modifie)
}
\end{lstlisting}

\subsection{Passage par référence}

Le paramètre est un \textbf{alias} de la variable originale.

\begin{lstlisting}[style=cpp, caption={Passage par référence}]
void modifier(int& x) {   // Le & est crucial
    x = 10;          // Modifie l'ORIGINAL
}

int main() {
    int a = 5;
    modifier(a);
    cout << a;       // Affiche : 10 (modifie !)
}
\end{lstlisting}

\subsection{Passage par référence constante}

Permet d'éviter la copie tout en protégeant la variable.

\begin{lstlisting}[style=cpp, caption={Passage par référence constante}]
void afficher(const int& x) {
    cout << x;
    // x = 10;      // ERREUR de compilation
}
\end{lstlisting}

\begin{regle}
Quand utiliser quoi ?
\begin{itemize}
    \item \textbf{Par valeur} : types simples (\texttt{int}, \texttt{char}, \texttt{bool}, \texttt{double})
    \item \textbf{Par référence constante} (\texttt{const \&}) : objets en lecture seule
    \item \textbf{Par référence} (\texttt{\&}) : objets à modifier
\end{itemize}
\end{regle}

\subsection{Tableau comparatif}

\begin{table}[h!]
\centering
\caption{Modes de passage de paramètres}
\begin{tabular}{llll}
\toprule
\textbf{Mode} & \textbf{Copie ?} & \textbf{Modifiable ?} & \textbf{Syntaxe} \\
\midrule
Par valeur & Oui & Non (copie) & \texttt{void f(int x)} \\
Par référence & Non & Oui & \texttt{void f(int\& x)} \\
Par réf. const. & Non & Non & \texttt{void f(const int\& x)} \\
Par pointeur (C) & Non & Oui & \texttt{void f(int* x)} \\
\bottomrule
\end{tabular}
\end{table}

% ============================================================
\section{Programmation orientée objet}

\subsection{Concepts fondamentaux}

La \textbf{Programmation Orientée Objet} (POO) permet de :
\begin{itemize}
    \item Regrouper données et fonctions dans une même structure (classe)
    \item Protéger les données (encapsulation)
    \item Créer des types personnalisés
    \item Améliorer la maintenabilité et la réutilisabilité du code
\end{itemize}

\begin{astuce}
Le lien avec les structures C est direct : une \textbf{classe} en C++ est comme une \texttt{struct} en C, mais avec des \textbf{méthodes intégrées} et un contrôle d'accès (\texttt{private} / \texttt{public}).
\end{astuce}

\subsection{Structure d'une classe}

\begin{lstlisting}[style=cpp, caption={Déclaration d'une classe - fichier Point.h}]
#ifndef POINT_H
#define POINT_H

class Point {
private:     // Donnees membres privees
    double x;
    double y;

public:      // Fonctions membres publiques
    // Constructeur
    Point(double abs, double ord);

    // Methodes
    void afficher() const;
    void deplacer(double dx, double dy);

    // Accesseurs (getters)
    double getX() const;
    double getY() const;

    // Mutateurs (setters)
    void setX(double valeur);
    void setY(double valeur);
};

#endif
\end{lstlisting}

\subsection{Implémentation}

\begin{lstlisting}[style=cpp, caption={Implémentation de la classe - fichier Point.cpp}]
#include "Point.h"
#include <iostream>
using namespace std;

// Constructeur
Point::Point(double abs, double ord) {
    x = abs;
    y = ord;
}

void Point::afficher() const {
    cout << "(" << x << ", " << y << ")" << endl;
}

void Point::deplacer(double dx, double dy) {
    x += dx;
    y += dy;
}

double Point::getX() const { return x; }
double Point::getY() const { return y; }
void Point::setX(double valeur) { x = valeur; }
void Point::setY(double valeur) { y = valeur; }
\end{lstlisting}

\subsection{Utilisation}

\begin{lstlisting}[style=cpp, caption={Programme principal - fichier main.cpp}]
#include "Point.h"
#include <iostream>
using namespace std;

int main() {
    Point p1(3.0, 5.0);
    Point p2(1.0, 2.0);

    p1.afficher();           // (3, 5)
    p1.deplacer(2.0, 1.0);
    p1.afficher();           // (5, 6)

    cout << "x = " << p1.getX() << endl;

    p1.setX(10.0);
    p1.afficher();           // (10, 6)

    // p1.x = 20;            // ERREUR : donnee privee

    return 0;
}
\end{lstlisting}

\begin{regle}
Principe d'encapsulation :
\begin{itemize}
    \item Les données membres sont \texttt{private}
    \item Les méthodes d'interface sont \texttt{public}
    \item On accède aux données via des accesseurs (getters) et mutateurs (setters)
\end{itemize}
\end{regle}

% ============================================================
\section{Constructeurs et destructeurs}

\subsection{Le constructeur}

Le constructeur est une fonction appelée \textbf{automatiquement} lors de la création d'un objet.

\begin{regle}
Caractéristiques :
\begin{itemize}
    \item Porte le même nom que la classe
    \item N'a pas de type de retour
    \item Sert à initialiser l'objet
    \item Peut effectuer des allocations dynamiques
\end{itemize}
\end{regle}

\subsection{Types de constructeurs}

\begin{lstlisting}[style=cpp, caption={Différents types de constructeurs}]
class Tableau {
private:
    int* tab;
    int taille;

public:
    // 1) Constructeur par defaut
    Tableau() {
        taille = 10;
        tab = new int[10];
    }

    // 2) Constructeur avec parametres
    Tableau(int n) {
        taille = n;
        tab = new int[n];
    }

    // 3) Constructeur avec valeur par defaut
    // Tableau(int n = 10) { taille = n; tab = new int[n]; }
};
\end{lstlisting}

\subsection{Utilisation}

\begin{lstlisting}[style=cpp, caption={Appel des constructeurs}]
int main() {
    Tableau t1;          // Appelle le constructeur par defaut
    Tableau t2(20);      // Appelle le constructeur avec n=20
    Tableau t3(50);      // Appelle le constructeur avec n=50

    // ATTENTION : ceci ne cree PAS un objet !
    // Tableau t4();     // Declare une fonction !
}
\end{lstlisting}

\subsection{Le destructeur}

Le destructeur est une fonction appelée \textbf{automatiquement} lors de la destruction d'un objet.

\begin{regle}
Caractéristiques :
\begin{itemize}
    \item Porte le nom de la classe précédé de \texttt{\~{}}
    \item N'a pas de paramètres ni de type de retour
    \item Sert à libérer les ressources (mémoire, fichiers, etc.)
\end{itemize}
\end{regle}

\begin{lstlisting}[style=cpp, caption={Destructeur}]
class Tableau {
private:
    int* tab;
    int taille;

public:
    Tableau(int n) {
        taille = n;
        tab = new int[n];
        cout << "Construction d'un tableau de " << n << endl;
    }

    ~Tableau() {
        delete[] tab;
        cout << "Destruction du tableau" << endl;
    }
};
\end{lstlisting}

% ============================================================
\section{Constructeur par recopie et opérateur d'affectation}

\subsection{Problème de la copie superficielle}

Quand une classe contient un pointeur, la copie par défaut ne copie que l'adresse — les deux objets pointent alors sur la même zone mémoire. Cela provoque des erreurs lors de la destruction.

\begin{lstlisting}[style=cpp, caption={Problème de copie superficielle}]
class Tableau {
    int* tab;
    int taille;
public:
    Tableau(int n) { taille = n; tab = new int[n]; }
    ~Tableau() { delete[] tab; }
};

int main() {
    Tableau t1(5);
    Tableau t2 = t1;    // Copie superficielle !
    // t1 et t2 partagent le MEME tableau
    // A la destruction : double delete -> crash !
}
\end{lstlisting}

\subsection{Constructeur par recopie}

\begin{lstlisting}[style=cpp, caption={Constructeur par recopie (copie profonde)}]
// Dans la classe :
Tableau(const Tableau& autre) {
    taille = autre.taille;
    tab = new int[taille];           // Nouvelle allocation
    for (int i = 0; i < taille; i++) {
        tab[i] = autre.tab[i];       // Copie des valeurs
    }
}
\end{lstlisting}

\subsection{Opérateur d'affectation}

\begin{lstlisting}[style=cpp, caption={Opérateur d'affectation}]
Tableau& operator=(const Tableau& autre) {
    if (this == &autre)              // Auto-affectation
        return *this;

    delete[] tab;                    // Libere l'ancienne memoire

    taille = autre.taille;
    tab = new int[taille];
    for (int i = 0; i < taille; i++) {
        tab[i] = autre.tab[i];
    }

    return *this;
}
\end{lstlisting}

\begin{regle}
\textbf{Règle des trois} : si vous définissez l'un des trois suivants, définissez les trois :
\begin{enumerate}
    \item Destructeur
    \item Constructeur par recopie
    \item Opérateur d'affectation
\end{enumerate}
\end{regle}

% ============================================================
\section{Listes chaînées}

\subsection{Principe}

Une liste chaînée est une structure de données où chaque élément (\textbf{maillon} ou \textbf{nœud}) contient une valeur et un \textbf{pointeur vers le maillon suivant}.

\begin{lstlisting}[style=cpp, caption={Structure d'un maillon}]
class Maillon {
public:
    int valeur;
    Maillon* suivant;

    Maillon(int v) : valeur(v), suivant(nullptr) {}
};
\end{lstlisting}

\subsection{Classe ListeChainee}

\begin{lstlisting}[style=cpp, caption={Liste chaînée complète}]
class ListeChainee {
private:
    Maillon* tete;

public:
    ListeChainee() : tete(nullptr) {}

    ~ListeChainee() {
        Maillon* courant = tete;
        while (courant != nullptr) {
            Maillon* suivant = courant->suivant;
            delete courant;
            courant = suivant;
        }
    }

    // Ajouter en tete
    void ajouterEnTete(int valeur) {
        Maillon* nouveau = new Maillon(valeur);
        nouveau->suivant = tete;
        tete = nouveau;
    }

    // Parcourir et afficher
    void afficher() const {
        Maillon* courant = tete;
        while (courant != nullptr) {
            cout << courant->valeur << " -> ";
            courant = courant->suivant;
        }
        cout << "NULL" << endl;
    }

    bool estVide() const { return tete == nullptr; }
};
\end{lstlisting}

\subsection{Comparaison tableau / liste chaînée}

\begin{table}[h!]
\centering
\caption{Tableau vs Liste chaînée}
\begin{tabular}{lll}
\toprule
 & \textbf{Tableau} & \textbf{Liste chaînée} \\
\midrule
Taille & Fixe & Dynamique \\
Accès à l'élément i & O(1) (direct) & O(n) (parcours) \\
Insertion en tête & O(n) & O(1) \\
Mémoire & Continue & Fragmentée \\
\bottomrule
\end{tabular}
\end{table}

% ============================================================
\section{Exercices récapitulatifs}

\subsection{Exercice 1 : Classe Rectangle}

\textbf{Énoncé :} Créer une classe \texttt{Rectangle} avec longueur et largeur, un constructeur, et les méthodes \texttt{aire()}, \texttt{perimetre()} et \texttt{afficher()}.

\textbf{Solution :}

\begin{lstlisting}[style=cpp, caption={Rectangle.h}]
#ifndef RECTANGLE_H
#define RECTANGLE_H

class Rectangle {
private:
    double longueur;
    double largeur;

public:
    Rectangle(double l, double L);
    double aire() const;
    double perimetre() const;
    void afficher() const;
};

#endif
\end{lstlisting}

\begin{lstlisting}[style=cpp, caption={Rectangle.cpp}]
#include "Rectangle.h"
#include <iostream>
using namespace std;

Rectangle::Rectangle(double l, double L) {
    longueur = l;
    largeur = L;
}

double Rectangle::aire() const {
    return longueur * largeur;
}

double Rectangle::perimetre() const {
    return 2 * (longueur + largeur);
}

void Rectangle::afficher() const {
    cout << "Rectangle " << longueur << "x" << largeur << endl;
    cout << "Aire : " << aire() << endl;
    cout << "Perimetre : " << perimetre() << endl;
}
\end{lstlisting}

\subsection{Exercice 2 : Classe Vecteur dynamique}

\textbf{Énoncé :} Créer une classe \texttt{Vecteur} qui gère un tableau dynamique d'entiers avec :
\begin{itemize}
    \item Constructeur avec taille
    \item Destructeur
    \item Constructeur par recopie
    \item Opérateur d'affectation
    \item Méthode \texttt{get(int i)} et \texttt{set(int i, int valeur)}
\end{itemize}

\textbf{Solution :}

\begin{lstlisting}[style=cpp, caption={Vecteur.h}]
#ifndef VECTEUR_H
#define VECTEUR_H

class Vecteur {
private:
    int* tab;
    int taille;

public:
    Vecteur(int n);
    ~Vecteur();
    Vecteur(const Vecteur& autre);
    Vecteur& operator=(const Vecteur& autre);

    int get(int i) const;
    void set(int i, int valeur);
    int getTaille() const;
    void afficher() const;
};

#endif
\end{lstlisting}

\begin{lstlisting}[style=cpp, caption={Vecteur.cpp}]
#include "Vecteur.h"
#include <iostream>
using namespace std;

Vecteur::Vecteur(int n) {
    taille = n;
    tab = new int[n];
    for (int i = 0; i < n; i++) tab[i] = 0;
}

Vecteur::~Vecteur() { delete[] tab; }

Vecteur::Vecteur(const Vecteur& autre) {
    taille = autre.taille;
    tab = new int[taille];
    for (int i = 0; i < taille; i++) tab[i] = autre.tab[i];
}

Vecteur& Vecteur::operator=(const Vecteur& autre) {
    if (this == &autre) return *this;
    delete[] tab;
    taille = autre.taille;
    tab = new int[taille];
    for (int i = 0; i < taille; i++) tab[i] = autre.tab[i];
    return *this;
}

int Vecteur::get(int i) const {
    if (i < 0 || i >= taille) {
        cout << "Erreur : indice hors limites" << endl;
        return 0;
    }
    return tab[i];
}

void Vecteur::set(int i, int valeur) {
    if (i < 0 || i >= taille) {
        cout << "Erreur : indice hors limites" << endl;
        return;
    }
    tab[i] = valeur;
}

int Vecteur::getTaille() const { return taille; }

void Vecteur::afficher() const {
    cout << "[";
    for (int i = 0; i < taille; i++) {
        cout << tab[i];
        if (i < taille - 1) cout << ", ";
    }
    cout << "]" << endl;
}
\end{lstlisting}

\subsection{Exercice 3 : Pile (Stack) avec liste chaînée}

\textbf{Énoncé :} Implémenter une pile (LIFO) avec :
\begin{itemize}
    \item \texttt{empiler(int x)} : ajoute x au sommet
    \item \texttt{depiler()} : retire et retourne l'élément du sommet
    \item \texttt{estVide()} : teste si la pile est vide
    \item \texttt{sommet()} : retourne l'élément du sommet sans le retirer
\end{itemize}

\textbf{Solution :}

\begin{lstlisting}[style=cpp, caption={Pile.h et Pile.cpp}]
class Maillon {
    int valeur;
    Maillon* suivant;
    friend class Pile;
public:
    Maillon(int v, Maillon* s = nullptr) : valeur(v), suivant(s) {}
};

class Pile {
private:
    Maillon* sommetPile;

public:
    Pile() : sommetPile(nullptr) {}

    ~Pile() {
        while (!estVide()) depiler();
    }

    void empiler(int x) {
        sommetPile = new Maillon(x, sommetPile);
    }

    int depiler() {
        if (estVide()) { cout << "Erreur : pile vide" << endl; return 0; }
        int valeur = sommetPile->valeur;
        Maillon* temp = sommetPile;
        sommetPile = sommetPile->suivant;
        delete temp;
        return valeur;
    }

    bool estVide() const { return sommetPile == nullptr; }

    int sommet() const {
        if (estVide()) { cout << "Erreur : pile vide" << endl; return 0; }
        return sommetPile->valeur;
    }
};
\end{lstlisting}

% ============================================================
\section{Aide-mémoire}

\subsection{Tableau récapitulatif des opérateurs}

\begin{table}[h!]
\centering
\caption{Opérateurs importants en C / C++}
\begin{tabular}{llll}
\toprule
\textbf{Opérateur} & \textbf{Nom} & \textbf{C} & \textbf{C++} \\
\midrule
\texttt{\&} & Adresse & \checkmark & \checkmark \\
\texttt{*} & Déréférencement / pointeur & \checkmark & \checkmark \\
\texttt{->} & Accès membre via pointeur & \checkmark & \checkmark \\
\texttt{.} & Accès membre d'un objet & \checkmark & \checkmark \\
\texttt{::} & Résolution de portée & \textemdash & \checkmark \\
\texttt{new} & Allocation dynamique & \textemdash & \checkmark \\
\texttt{delete} & Libération dynamique & \textemdash & \checkmark \\
\texttt{malloc} & Allocation dynamique & \checkmark & (\checkmark) \\
\texttt{free} & Libération dynamique & \checkmark & (\checkmark) \\
\bottomrule
\end{tabular}
\end{table}

\subsection{Commandes de compilation}

\begin{lstlisting}[style=bash, caption={Compilation C et C++}]
# Compilation C
gcc programme.c -o programme

# Compilation C++
g++ programme.cpp -o programme

# Avec debug
g++ -g programme.cpp -o programme

# Avec warnings
g++ -Wall -Wextra programme.cpp -o programme

# Plusieurs fichiers
g++ main.cpp Point.cpp Rectangle.cpp -o programme

# Execution
./programme

# Detection de fuites memoire (Linux)
valgrind ./programme
\end{lstlisting}

\subsection{Modèle de fichiers .h et .cpp (C++)}

\begin{lstlisting}[style=cpp, caption={Modèle de fichier .h}]
#ifndef NOMCLASSE_H
#define NOMCLASSE_H

class NomClasse {
private:
    // Attributs prives

public:
    NomClasse();                              // Constructeur
    NomClasse(/* parametres */);              // Constructeur parametré
    ~NomClasse();                             // Destructeur
    NomClasse(const NomClasse& autre);        // Recopie
    NomClasse& operator=(const NomClasse& autre); // Affectation

    // Methodes publiques
};

#endif
\end{lstlisting}

% ============================================================
\section{Conclusion}

Ce document couvre l'ensemble du parcours C/C++ du semestre 1 et 2 de licence informatique. En résumé :

\textbf{Partie C :}
\begin{itemize}
    \item Algorithmes, variables, conditions, boucles
    \item Fonctions : portée, passage par valeur, retour, fichiers .h/.c
    \item Pointeurs : \texttt{\&}, \texttt{*}, passage in-out
    \item Tableaux statiques et lien avec les pointeurs
    \item Structures : regrouper des données cohérentes
    \item Allocation dynamique : \texttt{malloc / free}, tas vs pile
\end{itemize}

\textbf{Partie C++ :}
\begin{itemize}
    \item Syntaxe C++ : \texttt{cout/cin}, \texttt{string}, références
    \item Allocation dynamique : \texttt{new / delete}
    \item Programmation orientée objet : classes, encapsulation
    \item Constructeurs et destructeurs
    \item Règle des trois : destructeur, recopie, affectation
    \item Listes chaînées
\end{itemize}

\begin{regle}
Les trois règles d'or du C++ :
\begin{enumerate}
    \item \textbf{Règle des trois} : Destructeur + constructeur par recopie + opérateur d'affectation
    \item \textbf{RAII} : Acquérir les ressources dans le constructeur, les libérer dans le destructeur
    \item \textbf{const-correctness} : Utiliser \texttt{const} partout où c'est pertinent
\end{enumerate}
\end{regle}

\begin{center}
\textbf{Bon courage pour les révisions !}
\end{center}

\end{document}